% gwebman-ko.tex -- the GWEB manual in Korean, a translation of gwebman.tex.
% Process with luatex (Hangul needs luatexko), for example
%       TEXINPUTS=..: luatex gwebman-ko.tex
% or simply `make manual-ko' from the top of the repository.
%
% It keeps gwebman.tex's own macros and layout, and loads kotexgweb.tex for the
% Hangul fonts. When gwebman.tex changes, change this translation with it.

\input gwebmac
\input kotexgweb
\pdffalse \dvipdffalse        % a plain manual: no hyperlinks or bookmarks
\parskip=0pt plus 1pt
\parindent=18pt
\abovedisplayskip=.5\abovedisplayskip  \belowdisplayskip=.5\belowdisplayskip
\abovedisplayshortskip=.5\abovedisplayshortskip
\belowdisplayshortskip=.5\belowdisplayshortskip

\font\eighttt=cmtt8
\font\ninett=cmtt9
\font\twelvett=cmtt12
% The 12-point Latin faces of gwebman.tex, each paired with a Hangul face of the
% same weight, the way kotexgweb pairs \tenbf with \tenbd. The Hangul face is
% set at 0.962 of the size, the ratio hangulfont.tex uses for the serif.
\font\latinauthorfont=cmr12
\font\latinchapterfont=cmbx12
\sethangulfont\twelvebd="Noto Serif CJK KR Bold:%
  language=KOR;script=hang;interhangul=-0.04em;interlatincjk=.125em;%
  charraise=-0.06em;+compresspunctuations;expansion=default;" at 11.54pt
\def\authorfont{\latinauthorfont\twelvemj}
\def\chapterfont{\latinchapterfont\twelvebd}

%---inline code; @ < > = etc. are safe in \tt, so avoid braces inside \cw ---
\def\cw#1{{\tt #1}}
\def\bs{\char92 }                              % a literal backslash, in \tt
\def\rn#1{$\langle\,${\it #1\/}$\,\rangle$}    % a refinement name, as in print

%---a centred chapter heading, in the manner of cwebman's \section ---
\def\chapter#1{\penalty-500\bigskip
  \centerline{\chapterfont\def\.##1{{\twelvett##1}}#1}\nobreak\vskip 6pt}
%---a run-in bold subsection heading ---
\def\section#1{\medskip\noindent{\bf #1.}\enspace\ignorespaces}

%---a bulleted list item ---
\def\bull{\smallskip\item{$\bullet$}}

%---a verbatim code display (TeXbook recipe): | is the escape inside it ---
\def\begincode{\par\medskip\begingroup
  \def\do##1{\catcode`##1=12 }\dospecials
  \catcode`\|=0
  \obeylines \obeyspaces \tt \parindent=0pt \leftskip=\parindent\relax}
\def\endcode{\endgroup\medskip\noindent}

%---running heads, in the manner of cwebman ---
\def\runninghead{{\ninett GWEB} 사용자 설명서 (버전 0.10.7)}
\def\lheader{\mainfont\folio\hfil{\eightrm\eightgt\runninghead}\hfil}
\def\rheader{\hfil{\eightrm\eightgt\runninghead}\hfil\mainfont\folio}

% =============================================================================
% Title page (no running head), in the manner of cwebman.
\headline={\hfil}\footline={\hfil}
\null\vfill
\centerline{\titlefont {\ttitlefont GWEB} 시스템}
\vskip 6pt
\centerline{\chapterfont Go를 위한 문학적 프로그래밍 시스템}
\vskip 18pt
\centerline{(버전 0.10.7)}
\vskip 24pt
\centerline{\authorfont 크누스와 레비의 {\ttitlefont CWEB}을 본떠 만들다}
\vfill
\noindent
이 글은 \.{CWEB}이나 문학적 프로그래밍을 이미 아는 독자를 위한 \.{GWEB}
안내서다. 문학적 프로그래밍이 무엇인지는 다시 설명하지 않는다. 그 대신
\.{GWEB}이 \.{CWEB}과 어떻게 다른지, 그리고 실제로 쓰게 될 기능과 명령에
집중한다. 대강 말해 \.{GWEB}은 ``\CEE/를 \GO/로 바꾼 \.{CWEB}''이므로,
\.{CWEB}에 관해 아는 것은 대부분 그대로 통한다. 흥미로운 대목은 \GO/와
\cw{gofmt}가 그림을 바꾸는 곳이다.
\pageno=0
\vfil\eject

% Running heads from the second page on.
\headline={\ifodd\pageno\rheader\else\lheader\fi}

\chapter{한눈에 보는 \.{GWEB}}
\noindent\.{GWEB}은 \.{CWEB}의 도구를 하나하나 그대로 본뜬다.

\bull \cw{gtangle} $\approx$ \cw{ctangle}---컴파일할 수 있는 \GO/ 소스를 뽑아낸다.

\bull \cw{gweave} $\approx$ \cw{cweave}---문서를 조판한다.

\bull \cw{gwebmac.tex} $\approx$ \cw{cwebmac.tex}---짜낸 출력이 쓰는 \TeX\ 매크로다.

\noindent 웹은 파일 하나이고, 관례상 이름을 \cw{foo.w}로 짓는다. 명령
\cw{gtangle foo}는 \cw{foo.go}를(덧붙인 출력 파일이 있으면 그것도) 쓰고, 명령
\cw{gweave foo}는 \cw{foo.tex}를 쓴다. 이 파일을 \cw{pdftex}로 돌리면
\cw{foo.pdf}가 나온다. \.{CWEB}처럼 명령줄에서는 \cw{.w}를 빼도 된다.

\chapter{\.{CWEB}에서 그대로 가져온 것}
\noindent
\.{CWEB}을 안다면 \.{GWEB}도 거의 다 아는 셈이다. 웹은 번호가 붙은
{\it 절\/}이 이어진 것이고, 각 절에는 차례대로 세 부분까지 올 수 있다. 바로
\TeX\ 설명 파트, 정의 파트(\cw{@d}/\cw{@f}/\cw{@s}), 코드 파트다. 그리고 이
절들 앞에는 {\it 림보\/}가 온다. 제어 코드는 익숙한 뜻을 그대로 지닌다.

\bull 제어 코드 \cw{@*}, \cw{@*1}, \cw{@*2}, $\ldots$는 별표 절(장)을 시작하고, \cw{@\ }는
보통 절을 시작한다.

\bull 제어 코드 \cw{@c}는 이름 없는 코드를 시작한다. 제어 코드 \cw{@<Name@>=}는 이름 있는 절을
정의하고, \cw{@<Name@>}는 그것을 참조한다. 이름 끝에 \cw{...}를 붙이면 그
이름을 유일하게 가리키는 앞부분만으로 줄여 쓸 수 있다.

\bull 제어 코드 \cw{@(file@>=}는 뒤따르는 코드를 다른 출력 파일로 보내고, \cw{@i}는
다른 웹을 끼워 넣는다.

\bull 지시자 \cw{@f}/\cw{@s}는 식별자의 조판 방식을 정한다. 색인 코드 \cw{@\char94 text@>},
\cw{@.text@>}, \cw{@:text@>}는 색인 항목을 더하고, \cw{@!}는 그 자리를
색인에서 정의로 표시하게 한다.

\bull 제어 코드 \cw{@=...@>}는 원문 그대로 내보내고, \cw{@t...@>}는 \TeX\ 글이다.
제어 코드 \cw{@\char38}는 토큰을 이어 붙이고, \cw{@@}는 글자 \cw{@} 그 자체이며,
\cw{@\char37}와 \cw{@q...@>}는 주석이다.

\bull 배치 코드 \cw{@,} \cw{@/} \cw{@|} \cw{@\#}는 짜낸 문서의 배치를 손본다.

\bull 변경 파일(\cw{@x}/\cw{@y}/\cw{@z})과 설명 속의 인라인 코드 \cw{|...|}는
\.{CWEB}과 똑같이 동작한다.

\noindent 작은 편의가 둘 있다. 제어 코드 \cw{@p}는 \.{CWEB}처럼 \cw{@c}와 같은
뜻으로 받아들인다. 그리고 \cw{\bs input gwebmac}는 직접 쓸 필요가 없다.
명령 \cw{gweave}가 대신 써 넣고, 직접 넣어 둔 것이 있으면 지운다.

\chapter{\.{GWEB}이 다른 점}
\noindent 차이는 모두 \CEE/ 대신 \GO/를 조판한다는 데서 나온다.

\section{전처리기가 없다}
\GO/에는 전처리기가 없으므로 \cw{@d}는 결코 코드를 만들지 않는다. 명령
\cw{gtangle}은 이것을 무시한다(진짜 정의에는 \GO/의 \cw{const}와 \cw{func}를
쓴다). 그렇다고 쓸모가 없지는 않다. 명령 \cw{gweave}는 \cw{@d}를 \.{CWEB}의
매크로처럼 타자체로 찍을 식별자의 목록으로 읽는다(아래 참고). 매크로 하나만
받는 \.{CWEB}의 \cw{@d}와 달리 이름을 얼마든지 받는다. 본문은 다음 \cw{@}까지
이어지고 여러 줄에 걸쳐도 되므로, \cw{@d Push Pop Peek} 한 줄이면 세 이름이
한꺼번에 정해진다. 조판 지시자 \cw{@f}와 \cw{@s}도 남아 있다. 이를테면
지시자 \cw{@f Stack int}는 사용자 형 \cw{Stack}을 미리 선언된 형처럼 굵게 찍는다.

\section{파서 없이 보기 좋게 조판하기}
명령 \cw{cweave}는 \CEE/를 파싱해 모든 공백과 줄바꿈과 들여쓰기를 문법에서
새로 끌어낸다. 명령 \cw{gweave}도 수식 같은 결과에 이른다. 곱셈
\cw{a * b}는 띄우고, 포인터 \cw{*T}는 붙이고, 관계 연산자는 떼어 놓고,
키워드 뒤에는 한 칸을 둔다. 다만 {\it 국소적\/} 규칙만으로 그렇게 한다. 각
토큰과 그 이웃만 볼 뿐, 구문 트리는 만들지 않는다. 그래서 간격은
\cw{gofmt}가 아니라 \cw{cweave}를 따른다. 명령 \cw{gofmt}는 우선순위가 높은
항을 붙여 쓰지만(\cw{a*b + c}), 명령 \cw{gweave}는 수식을 조판하듯 모든
연산자를 똑같이 띄운다. 들여쓰기도 마찬가지로 구조를 따른다({\it 짜낸
문서\/} 참고). 그래서 들쭉날쭉하거나 빽빽한 소스도 깔끔하게 놓인다. 절을
손으로 다듬을 필요가 없다.

국소적으로는 답이 없는 모호함이 두 가지 있다. 이때 \cw{gweave}는 저자가
믿을 만하게 주는 단 하나의 신호, 곧 저자 자신의 띄어쓰기에 기댄다.
명령 \cw{cweave}가 \cw{@[}\thinspace$\ldots$\thinspace\cw{@]}에 기대는 것과 똑같다.
피연산자 사이의 \cw{*}는 소스에서 뒤에 붙여 쓰면(\cw{p *T}) 포인터 형이고,
그렇지 않으면(\cw{a * b}, \cw{a*b}) 곱셈이다. 중괄호는 그 앞의 키워드나 형에
따라 띄어 쓰는 블록을 열거나, 붙여 쓰는 합성 리터럴을 연다. \.{CWEB}의
조판 힌트 \cw{@+}, \cw{@[}\thinspace$\ldots$\thinspace\cw{@]}, \cw{@;}는 여기서
할 일이 없어 무시한다. 명시적 배치 코드 \cw{@,} \cw{@/} \cw{@|} \cw{@\#}는
여전히 자동 간격보다 앞서므로, 손으로 마무리할 때 쓰면 된다.

\section{tangle한 코드는 gofmt를 거친다}
명령 \cw{gtangle}이 내놓는 출력은 \cw{gofmt}를 거치므로, 올바른 웹은 늘 표준 형식의 \GO/를
낳는다. 모아 붙인 글이 올바른 \GO/가 아니면, 명령 \cw{gtangle}은 실패하는
대신 형식을 다듬지 않은 출력을 남기고 경고를 찍는다.

\section{형과 상수는 저절로 분류된다}
명령 \cw{gweave}는 \cw{type} 선언을 찾아낸다. 한 줄짜리 꼴(\cw{type Name
...})과 괄호로 묶은 꼴을 모두 찾고, 선언된 이름은 어디에 나오든 미리 선언된
형처럼 {\bf 굵게} 찍는다. 이름에 소문자가 없는 상수, 이를테면 \cw{TRAP}이나
\cw{MAX\_N}은 어디에 나오든 타자체로 찍는다. 명령 \cw{cweave}가 대문자 이름을
찍는 방식과 같다. 다만 \cw{N}처럼 한 글자인 이름은 \cw{cweave}에서처럼
이탤릭으로 남는다. 그 밖의 상수는 한 줄짜리 \cw{const Name = ...}이든,
묶음이든, \cw{iota} 열거든 보통의 이탤릭 식별자다. 이 규칙은 상수에만
적용된다. 대문자로 쓴 변수나 형은 제 분류를 지킨다. 이름을 \cw{@d}에 적으면
\cw{cweave}가 \cw{@d} 매크로를 찍듯 타자체로 찍는다. 이 웹에서 선언하지
{\it 않은\/} 이름도 마찬가지다. 이를테면 \cw{@d http.StatusOK}는
\cw{StatusOK}를 어디서나 타자체로 찍는다. 자동 분류는 무엇이든
\cw{@f}/\cw{@s}로 바꿀 수 있다. 이를테면 \cw{@f Name int}는 \cw{Name}을 억지로
굵게 찍고, \cw{@f Name x}(아무 보통 식별자)는 평범한 이탤릭으로 되돌린다.

이런 지시자는 꾸러미나 리시버로 {\it 한정\/}할 수 있다. 지시자
\cw{@s foo.Bar int}는 \cw{foo.Bar}로 쓴 \cw{Bar}만 정한다. 그래서 우연히 이름이
같은 \cw{abc.Bar}(어떤 \cw{abc}의 필드 \cw{Bar})나 그냥 \cw{Bar}는 건드리지
않는다. 한정하지 않은 \cw{@s Bar int}는 모든 \cw{Bar}를 정한다. 한정자는
\cw{.} 앞에 쓴 식별자일 뿐이므로, 실제 꾸러미 이름이 아니어도 된다.

\section{\.{GWEB}은 \.{GWEB}으로 쓰였다}
\.{CWEB}처럼 \.{GWEB}도 문학적 프로그램이다. 각 도구는 제 \GO/ 옆에 둔
\cw{.w} 웹에서 tangle된다. 원본은 \cw{common/common.w}(함께 쓰는 파서),
\cw{cmd/gtangle/gtangle.w}(명령과 tangle 엔진),
\cw{cmd/gweave/gweave.w}(명령과 weave 엔진)이다. 그 옆의 \GO/는 이들에서
tangle한 것이고 바이트 단위까지 똑같이 재현된다. 이 고정점은 바꿀 때마다
\cw{make bootstrap}이 확인하므로, 읽는 프로그램과 돌아가는 프로그램이 서로
어긋날 수 없다. \.{GWEB}은 라이브러리가 아니라 명령들의 묶음이다. 그래서 각
엔진은 따로 가져다 쓸 수 있는 꾸러미가 아니라 제 명령의 \cw{main} 꾸러미 안에
산다.
\begincode
make tangle       # re-tangle the .w sources back into the Go tree
make bootstrap    # tangle into a scratch tree and verify it is byte-for-byte
                  # identical to the committed sources (the fixpoint)
|endcode
\noindent 작업 순서는 이렇다. 파일 \cw{.w}를 고치고, \cw{make tangle}을 돌린 뒤,
둘을 함께 커밋한다. 테스트도 \cw{.w} 소스 안에 있다(웹마다 {\it Tests\/}
장이 있다). 그래서 테스트는 나머지 생성된 \GO/와 함께 tangle되고 git에서
무시된다. 커밋되는 \cw{\_test.go}는 하나도 없다.

저장소 맨 위의 \cw{gweb.w}는 위의 세 웹을 읽는 차례대로 \cw{@i}로 끼워 넣는
주 웹이다. 그래서 명령 \cw{gweave}는 시스템 전체를 한 문서로 조판할 수 있다.
\begincode
make selfdoc      # -> build/gweb.pdf, GWEB woven as a literate program
|endcode
\noindent 그 결과는 \.{GWEB}이 스스로를 기록한 문서다. 다른 프로그램에 만들어
주는 것과 똑같이 보기 좋게 조판한 코드와 색인과 상호 참조가 들어 있다. 각 웹을
따로 짤 수도 있다. 이를테면 \cw{gweave cmd/gweave/gweave.w}는 제 제목 쪽과
색인을 갖춘 \cw{gweave}의 문서를 만든다. \.{CWEB}의 \cw{common.w},
\cw{ctangle.w}, \cw{cweave.w}가 그렇듯이 말이다. 합본 문서는 \cw{gweb.ch}를
적용한다. 이 변경 파일은 각 웹의 제목과 \cw{@** Index.} 절을 걷어내어,
\cw{gweb.w}의 것만 남긴다.

\chapter{짜낸 문서}
\noindent
명령 \cw{gweave}는 \cw{cweave}와 같은 문서 구성을 만든다. 절 머리는 굵은
글씨로 본문 앞에 붙고(큰 별표 절 앞에서는 쪽이 바뀐다), 두 단으로 짠 색인,
{\it 절 이름 목록\/}, 목차 쪽, 절 번호를 단 머리글이 있다. 한 절은 설명과
코드를 통틀어 한 쪽에 들어가면 한 쪽에 함께 놓인다. 코드 파트는 그 자체가 한
쪽보다 길지 않은 한 쪽 경계에서 나뉘지 않는다.

\section{코드를 찍는 방식}
예약어와 사용자 정의 형은 굵게, 그 밖의 식별자는 이탤릭으로, 문자열과 룬
리터럴은 타자체(\.{CWEB}처럼 \cw{cmtex10})로, 수와 주석은 로만체로 찍는다.
\GO/에 맞춘 손질도 몇 가지 있다. 시프트 연산자 \cw{<<}와 \cw{>>}는 붙은 기호
$\ll$과 $\gg$로 찍는다. 함수 이름이나 키워드 뒤에 \cw{(}가 오면 사이에 가는
공백을 두어 \cw{f (x)}가 엉겨 붙지 않게 한다. 그리고 주석의 \cw{//}는 조금 더
촘촘하게 붙인다.

코드 주석의 글은 \.{CWEB}과 똑같이 \TeX\ 그대로 옮긴다. 달러 기호 사이의
\TeX\ 수식은 수식으로 조판되고, 인라인 코드 \cw{|...|}는 그것이 가리키는
\GO/로 찍히며, \cw{\bs.\char123...\char125}는 타자체가 되고, 보통의 제어
시퀀스도 동작한다. 대가도 \.{CWEB}과 같다. \TeX\ 특수 문자가 살아 있으므로,
주석에 쓰는 글자 \cw{\char123}, \cw{\$}, \cw{\%}, \cw{\char38}, 역빗금은
막거나 짝을 맞춰야 한다(\cw{\bs\char123}, \cw{\bs\$}로 쓰고, 막대 자체는
\cw{\bs\char124}로 쓴다). 또 홀로 남은 \cw{\%}는 그 줄의 나머지를 주석으로
만들어 버린다. 이런 것이 많이 필요한 글은 코드가 아니라 \TeX\ 파트에 둔다.

들여쓰기도 간격도 소스에서 베끼지 않는다. 들여쓰기는 {\it 구조적\/}이다.
명령 \cw{gweave}는 각 코드 파트의 블록 중첩을 따라간다. 중괄호와 괄호와
대괄호를 세고, \cw{switch}의 \cw{case} 몸체와 레이블에 대해서는 \cw{gofmt}가
두는 특례를 따르며, 그에 맞춰 각 줄을 놓는다. 가로 간격도 마찬가지로 문법에서
끌어내며 수식과 같다({\it 파서 없이 보기 좋게 조판하기\/} 참고). 그래서
들여쓰기를 아무렇게나 했거나 공백 없이 빽빽하게 친 절도 깔끔하게 읽힌다.
명령 \cw{cweave}가 \CEE/ 문법에 따라 들여쓰기와 간격을 새로 정하는 것과 거의
같다.

\section{색인과 조판 코드}
손으로 다는 색인 코드 \cw{@\char94 text@>}(로만체), \cw{@.text@>}(타자체),
\cw{@:text@>}(\cw{\bs9}로 형식을 정함)는 \.{CWEB}처럼 동작한다. 정의 항목을
강제하는 \cw{@!}와, 이름의 분류를 바꾸는 \cw{@f}/\cw{@s}(\.{CWEB}처럼 같은 뜻)도
마찬가지다. 식별자가 \cw{func}, \cw{var}, \cw{const}, \cw{type} 뒤에 오면, 홀로
쓰였든 괄호 묶음의 한 항목으로 쓰였든 ``정의''로 표시한다(색인의 절 번호에
밑줄을 긋는다). 바로 오른쪽에 \cw{:=}가 오는 식별자도 그렇다. 구조체 \cw{struct}의
필드와 메서드의 이름도 \.{CWEB}에서 구조체 멤버가 그렇듯 정의로 친다. 다만
임베디드 필드는 형의 이름이므로 정의가 아니다. 이것은 형 분석이 아니라
어림셈이다. 한 번에 코드 파트 하나만 보므로, 다른 절에서 \cw{struct}에
끼워 넣은 필드는 알아채지 못한다. 그런 필드는 \cw{@!}로 표시하고, 나머지도
\cw{@!}로 바로잡는다. 형의 색인 항목은 그 형이 나오는 모든 곳과 마찬가지로
굵게 찍는다. \.{CWEB}이 \cw{typedef}를 찍는 방식과 같다.

\section{변경 표시} 변경 파일을 써서 짜면, 명령 \cw{gweave}는 그 변경이 손댄
모든 절의 번호 바로 뒤에 \cw{*}를 붙인다. 변경이 절의 줄을 고쳐 썼든 절을 새로
보탰든 마찬가지이고, \cw{cweave}와 똑같다. 이 표시는 그 번호가 나오는 곳마다
따라붙는다. 절 머리에도, 색인에도, 상호 참조 메모에도 붙는다. 색인 바로 앞의
작은 글씨 한 줄은 바뀐 절을 모두 나열한다. 색인 절 자체에도 표시가 붙는데,
어디가 바뀌든 색인에 찍히는 번호가 달라질 수 있기 때문이다. 림보에만 그친
변경은 제 절이 없으므로 아무것도 표시하지 않는다.

\section{하이퍼링크와 북마크}
참조로 찍힌 절 번호는 색인에서든, 상호 참조 메모에서든, 목차 쪽에서든 링크가
되고, 별표 절은 PDF 개요를 이룬다. \.{CWEB}과 똑같이 두 가지 백엔드가 이것을
만든다.

\bull PDF 모드에서 쓰는 엔진의 PDF 원시 명령. 엔진 \cw{pdftex}에는 이것이 바로 있다.
엔진 \cw{luatex}에는 \cw{\bs pdf} 원시 명령이 하나도 없다. 그것들이 모두
\cw{\bs pdfextension}의 다른 철자가 되었기 때문이다. 그래서 파일 \cw{gwebmac.tex}는
자신이 \cw{luatex}의 PDF 실행 위에 있음을 알아채면 그 이름들을
\cw{\bs pdfextension}에 이어 주고, 같은 코드 경로가 두 엔진을 다 섬긴다.
엔진 \cw{luatex}으로 짜야 하는 웹, 이를테면 \cw{luamplib}으로 MetaPost 그림을 그리는
웹이나 한글 웹도 그래서 따로 부탁하지 않아도 \cw{pdftex} 실행과 똑같은 링크와
개요를 얻는다.

\bull DVI를 거치는 길. \cw{\bs let\bs pdf+}로 \cw{\bs special} 백엔드를 요청하고,
그 결과를 \cw{dvipdfmx}로 바꾼다.

\noindent 매크로 \cw{\bs pdf}를 정의하지 않고 그냥 DVI로 짜면 링크와 개요는 빠진다.

\section{설명에 쓰는 도우미 매크로}
파일 \cw{gwebmac.tex}에는 서술에 쓰는 cwebmac의 편의 매크로가 들어 있다.
매크로 \cw{\bs.\char123...\char125}는 타자체 글을 찍고(특수 문자 자체는
\cw{\bs\bs}, \cw{\bs\char123}, $\ldots$로 쓴다), \cw{\bs\bs\char123 id\char125}는
이탤릭 식별자를 찍는다. 언어 이름표 \cw{\bs CEE/}, \cw{\bs GO/},
\cw{\bs UNIX/}, \cw{\bs TEX/}도 쓸 수 있다. 매크로 \cw{\bs,}는 수식에서든
글에서든 가는 공백이다. 그래도 \.{GWEB}에서 인라인 코드를 쓰는 관용구는 여전히
그냥 \cw{|...|}다.

\section{쪽 꾸미기}
다음을 림보에서 다시 정의한다(자동으로 들어가는 \cw{\bs input gwebmac} 뒤에서).

\bull \cw{\bs def\bs title\char123...\char125}---머리글과 목차의 제목이다.
기본값은 대문자로 바꾼 파일 이름이고, 머리글에는 모두 대문자로 나온다.

\bull \cw{\bs pagewidth}, \cw{\bs pageheight}---쪽 크기다.

\bull \cw{\bs def\bs contentspagenumber\char123 n\char125}와
\cw{\bs pageno}---더 큰 문서에 묶어 넣을 때 쪽 번호를 다시 매긴다.

\bull 레지스터 \cw{\bs codepenalty}(코드 파트가 쪽 나눔에 얼마나 세게 버티는지)와
\cw{\bs commentkern}(\cw{//}를 좁히는 정도)은 조정할 수 있는 레지스터다.

\bull \cw{\bs K}---\GO/의 짧은 선언 \cw{:=}의 기호다. 기본값은 두 글자를 붙여
찍는 것이다. \.{WEB}은 Pascal의 \cw{:=}를 왼쪽 화살표로 찍었는데, 림보에
\cw{\bs let\bs K=\bs Leftarrow}를 쓰면 여기서도 그렇게 된다.

\bull \cw{\bs AN}---\GO/의 비트 지우기 \cw{\&\string^}의 기호이고, 기본값은
동그라미 친 빗금이다. 기호 \cw{:=}처럼 본뜰 \.{CWEB}의 짝이 없으므로, 이것도
원하는 대로 \cw{\bs let}해 쓰면 된다.

\noindent 또 림보에 \cw{\bs datethis}를 쓰면 \.{CWEB}처럼 1절 아래에 실행한 날짜와
시각을 찍는다.

\.{CWEB}처럼 문서의 꼬리를 뺄 수도 있다. 림보에 다음 가운데 하나를 쓴다.
매크로 \cw{\bs nocon}은 목차를 빼고, \cw{\bs nosecs}는 절 이름 목록을(그리고
목차도) 빼고, \cw{\bs noinx}는 색인까지 뺀다(곧 프로그램 뒤의 모든 것을 뺀다).
매크로 \cw{\bs noatl}은 \.{CWEB}과 맞추려고 받아들이지만 아무 일도 하지 않는다.
\.{GWEB}에는 \cw{@l} 음역이 없기 때문이다.

목차 자료는 \cw{\bs contentsfile}에 쓰이고 \cw{\bs readcontents}가 다시 읽는다.
기본값은 각각 \cw{\bs jobname .toc}와 \cw{\bs input \bs contentsfile}이다.
림보에서 이 둘 가운데 하나를 다시 정의하면, \.{CWEB}에서처럼 따로 조판한 여러
조각의 목차를 하나로 모을 수 있다.

그리고 \cw{\bs pdfURL\char123 text\char125\char123 url\char125}를 \TeX\ 파트
어디에 쓰든, {\it text\/}를 파란 글씨로 찍고 {\it url\/}로 가는 링크로 만든다.
\.{CWEB}과 같다. PDF 백엔드가 없으면 둘 다 찍고 URL은 괄호 안에 넣는다. URL에서
at 기호는 \cw{@@}로 쓰고(\cw{.w}의 보통 탈출 방식이다), \cw{\char126}는
\cw{\bs TILDE/}로, \cw{\char95}는 \cw{\bs UNDER/}로 쓴다.

\chapter{다른 언어로 짜기}
\noindent 명령 \cw{gweave}가 짜낸 문서에 써 넣는 고정된 문구는 모두 매크로를
거친다. 코드 파트 아래의 상호 참조 메모, {\it 절 이름 목록\/}의 머리, 목차
쪽의 \cw{Section}과 \cw{Page}라는 말이 그렇다(\cw{\bs U}, \cw{\bs A},
\cw{\bs Nused}, \cw{\bs sectionword}, \cw{\bs pageword}, $\ldots$). 영어 글자를
그대로 쓰는 일은 없다. 이 매크로들을 \cw{\bs input}할 파일에서 다시 정의하면
문서를 어느 언어로든 옮길 수 있다. 그 파일을 유니코드를 다루는 엔진과 알맞은
글꼴과 짝지으면, 코드와 서술 자체도 ASCII가 아니어도 된다.

한글 지원은 이미 갖춰져 있다. 림보에 한 줄을 넣고
\begincode
\input kotexgweb
|endcode
평소처럼 짠 다음, 결과를 \cw{pdftex}가 아니라 \cw{luatex}으로 돌린다.
\begincode
gweave foo.w
luatex foo.tex
|endcode
파일 \cw{kotexgweb.tex}는 \cw{luatexko}와 Noto Serif/Sans CJK\thinspace KR
글꼴을 불러들이고, \cw{gweave}의 문구를 한국어로 옮긴다. 파란 상호 참조
링크와 북마크 개요는 \cw{gwebmac} 자체 덕분에 \cw{luatex}에서도 이미
동작한다(위의 {\it 하이퍼링크와 북마크\/} 참고). 파일 \cw{kotexgweb.tex}가
더하는 것은 한글 개요 제목에 필요한 UTF-16BE 변환뿐이다. 명령 \cw{gweave}
자체에는 특별한 플래그가 필요 없다. 글꼴은 파일에 심지 않고 이름으로 고르므로
미리 설치해 두어야 한다. 글꼴을 고르는 \cw{\bs sethangulfont} 줄들은 따로 떨어진
\cw{hangulfont.tex}(\cw{kotexgweb.tex}가 불러들인다)에 있으니, 다른 서체를
쓰고 싶으면 그 파일 하나만 고치면 된다. 예제 \cw{rps}는 이렇게 만든 완전한
웹이다. 그리고 다른 언어를 지원하려면 \cw{kotexgweb.tex}를 본보기로 삼아 베끼면
된다.

\chapter{프로그램 실행하기}
\begincode
gtangle [-o dir] file[.w] [change[.ch]]
gweave  [-o dir] file[.w] [change[.ch]]
|endcode
확장자 \cw{.w}는, 그리고 변경 파일의 \cw{.ch}도 빼도 된다. 옵션
\cw{-o}는 출력 디렉터리를 고르고, \cw{-version}은 버전을 찍고 끝낸다. 각
도구는 실행하면서 한 줄짜리 머리말을 찍고, 이어서 별표 절마다
\cw{*}{\it n\/} 진행 표시를 찍는다. \.{CWEB}의 도구와 같다.

명령 \cw{gtangle}은 tangle한 \GO/에 언제나 \cw{//line} 지시문을 넣는다.
\.{CWEB}의 \cw{ctangle}이 언제나 \cw{\#line}을 넣는 것과 같다. 그래서 컴파일러와
\cw{go vet}과 패닉 추적은 오류를 \cw{.go}가 아니라 \cw{.w}의 위치로 알려 준다.
\.{GWEB}은 제 소스도 이렇게 tangle하므로, \cw{.w}를 고치면 커밋된 \GO/의 그
줄 번호들이 뒤섞인다.

PDF를 만들려면 \cw{pdftex}로 곧장 가거나, 같은 하이퍼링크를 얻으려면 DVI와
\cw{dvipdfmx}를 거친다.
\begincode
pdftex foo.tex                       # native PDF, with links
tex "\let\pdf+ \input foo.tex"       # or via DVI...
dvipdfmx foo.dvi                     # ...then to PDF, with links
|endcode
그리고 \TeX\ 엔진이 \cw{gwebmac.tex}를 찾을 수 있어야 한다(아래의 설치 스크립트가 검색
경로에 넣어 준다. 그렇지 않으면 \cw{TEXINPUTS}를 정한다).

두 도구는 끝나지 않은 제어 코드, 정의되지 않은 절의 참조, 모호한 \cw{...}
줄임, 정의했지만 쓰이지 않은 절을 각각 \cw{file:line} 위치와 함께 알린다.
이것들은 경고다. 명령 \cw{gtangle}은 정의된 적 없는 참조를 실제로 펼쳐야 할
때에만 멈춘다.

\chapter{\.{GWEB} 설치하기}
\noindent \GO/로 커밋된 것은 \cw{gtangle}뿐이므로, 이것만 \GO/의 방식으로 설치할
수 있다.
\begincode
go install github.com/sjnam/gweb/cmd/gtangle@latest
|endcode
명령 \cw{gweave}는 커밋되지 않고 \cw{.w} 소스에서 생성되므로, 두 명령은
저장소 사본에서 함께 빌드한다. 명령 \cw{make install-tools}는 두 실행 파일을
\cw{GOBIN}에 넣는다.

완전히 설치하려면, 곧 명령들과, \cw{pdftex foo.tex}가 그냥 동작하도록 \TeX\
트리에 넣은 \cw{gwebmac.tex}와, 맨 페이지까지 설치하려면, 저장소 사본에서
\cw{./install.sh}를 돌린다(\cw{./install.sh --uninstall}은 이것을 지우고,
\cw{make install}은 이 스크립트를 부른다). 이 설명서는 \cw{make manual-ko}로
만들고, 영어 원본은 \cw{make manual}로 만든다.

\chapter{편집기 지원}
\noindent 디렉터리 \cw{editors/vscode} 아래의 VS\thinspace Code 확장이 \cw{.w} 파일을
지원한다. 구문 강조(제어 코드, 코드 파트와 \cw{|}$\ldots$\cw{|} 스팬에 들어 있는
\GO/, \TeX\ 설명)를 하고, {\bf 코드 파트 안에서 정의로 가기와 마우스 올리기}를
제공한다. 이 둘은 \cw{gtangle}이 출력에 남기는 \cw{//line} 소스 맵을 거쳐
\GO/ 확장으로 넘겨지므로, 한 웹에서 다른 웹에 정의된 코드로 건너뛰면 그 웹의
\cw{.w}에 닿는다. 여기에 \cw{@<section name@>} 이동과 자동 완성, 저장할 때 다시
tangle하기가 더해진다.

쓰려면 \pdfURL{릴리즈 페이지}{https://github.com/sjnam/gweb/releases}에서
\cw{.vsix}를 받아 설치하고(확장 보기 $\to$ {\it Install from VSIX\dots\/}),
편집기를 완전히 다시 띄운다. 자세한 것은 \cw{editors/vscode/README.md}를 본다.

\chapter{딸린 예제}
디렉터리 \cw{examples/}에는 간단한 것부터 풍성한 것까지 완성된 웹이 여덟 개
있다. 명령 \cw{make example}은 이것들을 모두 tangle하고 짜며, \cw{make -C
examples NAME=wc}는 하나만 처리한다.

\bull \cw{floyd.w}---플로이드의 분할 문제. 크누스가 {\it Are Toy Problems
Useful?}에서 다룬 고전적인 ``장난감 문제''로, $\sqrt1,\ldots,\sqrt{50}$을 거의
같은 두 반으로 나눈다. 가운데서 만나기 탐색, 그레이 코드 열거, 보정 합산,
\cw{math/big} 검증을 갖춘 문학적 풀이다.

\bull \cw{pairsums.w}---HackerRank의 {\it Pair Sums\/}. 모든 부분 배열에 걸친
쌍 곱의 합 가운데 가장 큰 것을 구한다. 항등식 ${\rm value} = (S^2-Q)/2$과
누적 합을 비튼 기법이 문제를 직선 무리의 위 포락선으로 바꾸고, 이것을
{\sl Li Chao 트리\/}로 $O(n\log n)$에 푼다.

\bull \cw{rps.w}---커피 내기로 위장한 한국어 문학적 수필. $N$명이 가위바위보로
누가 낼지 정하려면 몇 판이 걸리고, 가장 빨리 끝내려면 어떻게 무리를 나눠야
할까? 정확한 유리수로 \cw{math/big}을 써서 풀어 나간다. 기댓값 점화식, 정수
분할 전수 탐색, 100명을 위한 재귀 전략, 네 명씩 묶는 것이 최적이라는 증명,
볼록성 확인이 들어 있고, 짝이 되는 \cw{rps\_test.go}도 있다. 엔진 \cw{luatex}과
\cw{kotexgweb.tex}로 짠다({\it 다른 언어로 짜기\/} 참고).

\bull \cw{seq.w}---아주 작은 지연 수열 라이브러리(무한한 피보나치 수 위의
\cw{Map}, \cw{Filter}, \cw{Take}). \CEE/에는 짝이 없는 \GO/의 기능을 뽐낸다.
일급 함수와 클로저, 익명 함수, 제네릭, \GO/ 1.23의 range-over-func
반복자다.

\bull \cw{sham.w}---크누스의 Stanford GraphBase 시연 프로그램 \cw{sham}을
\.{GWEB}으로 옮긴 것. $8\times9$ 판에서 나이트 그래프의 대칭 해밀턴 회로를
센다. 그래프를 반으로 접고 \cw{goto} 레이블로 되추적한다. SGB를 \GO/로 옮긴
\pdfURL{go-sgb}{https://github.com/sjnam/go-sgb} 위에 지었으므로 돌리려면 그
모듈이 있어야 한다(\cw{go get github.com/sjnam/go-sgb}). 설명은 새로 썼다.
\.{GWEB}이 바깥 의존성과 진짜 크누스 프로그램을 다루는 모습을 보여 준다.

\bull \cw{squint.w}---요구에 따라 움직이는 채널 연결망으로 만든 지연 멱급수(합,
곱, 합성, 역수, 함수의 역, 그리고 \cw{exp} 같은 미분 방정식). 맥일로이의
{\it Squinting at Power Series\/}를 따랐다.

\bull \cw{torture.w}---일부러 엉망으로 만들었지만 컴파일은 되는, 1부터
번호를 매긴 펜윅 트리. 들여쓰기는 제멋대로이고 모든 줄의 코드가 빽빽하다.
그래서 이것을 짜 보면 문법을 따르는 조판기가 엉성한 소스를 얼마나 깔끔하게
다듬는지 드러난다.

\bull \cw{wc.w}---tangle한 출력이 시스템의 \cw{wc}와 똑같이 동작하는 낱말 세기
프로그램. 지시자 \cw{@f}와, CSV 변형을 만드는 변경 파일(\cw{wc.ch})을 보여 준다.

\noindent 명령 \cw{make test}(\cw{-short}가 아닌 실행)는 모든 예제를 tangle하고
그 결과를 \cw{go build}하므로, 예제가 언제나 컴파일되는 상태로 남는다. 더 많은
\.{GWEB} 예제는 \pdfURL{gweb-examples}{https://github.com/sjnam/gweb-examples}에
있다.

\noindent \.{GWEB}은 스스로를 기록하기도 한다. 명령 \cw{make selfdoc}은 시스템
자신의 \cw{.w} 소스를 PDF 하나로 짜고, \cw{make manual-ko}는 이 문서를 만든다.

\vfill
\centerline{\eightrm\eightmj {\eighttt GWEB}---Go를 위한 문학적 프로그래밍 시스템.}
\bye
